% ID: SEL-D0023-OB-E02
% ESTADO: cotejada. Segunda pasada visual completa y borrador compilado revisados.
% Fuente: Selectividad/matematicas_aplicadas/reserva_b_EXAMEN 2 -ANDALUCÍA.pdf
% SHA-256: 8cdb86259405652622f2ea129bffbdb239076ac612fd7fcebb86b9d5bdc451ec
% Página: 2; opción B; ejercicio original 2.

La función de costes de una fábrica, $f(x)$, en miles de euros, viene dada por la expresión:
\[
f(x)=2x^2-36x+200,
\]
donde $x$ es la cantidad fabricada del producto, en miles de kilogramos.
\begin{enumerate}[label=\alph*)]
\item \textbf{(0.8 puntos)} Determine la cantidad a fabricar para minimizar el coste y calcule este coste mínimo.
\item \textbf{(0.8 puntos)} A partir del signo de $f'(7)$, \textquestiondown{}qué se puede decir del coste para una producción de siete mil kilogramos?
\item \textbf{(0.9 puntos)} Dibuje la gráfica de la función de costes. \textquestiondown{}Para qué cantidad o cantidades fabricadas el coste es de 200000 \texteuro{}?
\end{enumerate}
